%===============================================================================
%      L'affichage
%===============================================================================
\section{Introduction}

   Testons \href{http://www.overleaf.com}{Something 
     Linky}
   
   Une des premières choses à faire, afin de nous simplifier la vie
pour la suite du développement, est de se doter d'un service minimal
d'affichage. Nous allons donc mettre en place un outil très simple
fondé sur ce que nous founit le matériel.

   Nous allons pour cela procéder en trois étapes.

\begin{itemize}
   \item   Dans un premier temps, nous allons écrire le code
permettant un accès simple à l'écran et une gestion minimale de ce
dernier dans un module ``{\em Console}''. Cet outil est suffisant pour
que notre noyau soit capable de nous donner un signe de vie à
l'écran. Mais nous n'allons pas nous arrêter en si bon chemin !

   \item Nous pourrons alors écrire une fonction \lstinline!printk!
utilisable partout dans le noyau et chargée de mettre en forme des
messages riches et de les envoyer dans le journal.

   \item   Nous construirons là-dessus un ``journal'' du noyau, c'est-à-dire
un outil permettant de rassembler et d'organiser les messages émis par
les différents sous-systèmes du noyau. Ils pourront être affichés sur
une console, ou conservés pour être stockés/traités d'une autre façon.

\end{itemize}

   L'avantage d'une telle structure est que nous pourrons la faire
évoluer. Nous pourrons par exemple créér des consoles virtuelles
fournissant le même service que notre outil de gestion de l'écran. Il
sera très simple de les intégrer dans le noyau.

%-------------------------------------------------------------------------------
%
%-------------------------------------------------------------------------------
\section{Gestion de l'écran : la console}

   Il se trouve que l'écran d'un PC peut être utilisé en mode textuel
comme un simple tableau de caractères. Ce tableau est situé à une
adresse bien connue $a$ et contient $l$ lignes de $c$
caractères. Chaque caractère est codé sur deux octets : l'un définit
le code {\sc ascii} du caractère affiché et l'autre la couleur du
texte et du fond. 

   Les fichiers \lstinline!manux/console.h! et \lstinline!lib/console.c!
implantent un outil de gestion de base de la console.

   On construit ici en particulier la fonction suivante, chargée
d'afficher sur l'écran une chaîne de caractères passée en paramètre :

\begin{lstlisting}
   void afficherEcran(char * msg);
\end{lstlisting}

   Elle prend en paramètre une chaîne de caractères ``classique''
(terminée par un caractère nul) et l'affiche à l'écran sans aucune
mise en forme.

   La gestion de la console est une des toutes premières choses
initialisées par le noyau, de sorte à pouvoir faire coucou dès que
possible.

%-------------------------------------------------------------------------------
%
%-------------------------------------------------------------------------------
\section{Un premier {\tt printk()}}

   Grâce à notre console , l'écriture de {\tt printk()} est assez
simple, puisque cette fonction n'a plus qu'à se préoccuper de créer
une chaîne de caractères et de la transmettre au journal (ou
directement à la console si on n'utilise pas le journal).

   Elle est mise en \oe{}uvre dans les fichiers
\lstinline!manux/printk.h! et \lstinline!noyau/printk.c!

   La fonction essentielle est bien sûr :

\begin{lstlisting}
   void printk(char * format, ...);
\end{lstlisting} 

   Dont le rôle est de mettre en forme une chaîne de caractères
qu'elle envoie ensuite à la console ou au journal.

%-------------------------------------------------------------------------------
%
%-------------------------------------------------------------------------------
\section{Plus de souplesse avec les consoles virtuelles}

   Lorsque nous allons mettre en place la notion de tâches, il sera
intéressant d'attribuer à chacune d'entre elles une console virtuelle
potentiellement différente des autres. Nous pourrons ainsi basculer de
l'une à l'autre pour mieux distinguer ce qu'affiche chaque tâche.

%-------------------------------------------------------------------------------
%
%-------------------------------------------------------------------------------
\section{Mise en place d'un journal}

   Nous allons définir un mécanisme de journal permettant de conserver
les messages transmis par le noyau. Pour le moment, notre journal se
contente de transmettre ces messages sur l'écran physique. Lorsque
nous aurons mis en place des outils de gestion mémoire, nous pourrons
étoffer ce mécanisme de journal.

   Le journal est implanté dans les fichiers
\lstinline!manux/journal.h! et \lstinline!lib/journal.c!. Il
propose essentiellement la fonction suivante

\begin{lstlisting}
   void journaliser(char * message); 
\end{lstlisting}

   Elle utilise la fonction \lstinline!afficherEcran()! pour envoyer le message à
l'écran.

   Évidemment, le journal est initialisé juste après la console.

%-------------------------------------------------------------------------------
%
%-------------------------------------------------------------------------------
\subsection{Introduction de plusieurs niveaux de messages}
   
%-------------------------------------------------------------------------------
%
%-------------------------------------------------------------------------------
\subsection{Utilisation des fichiers}

